\chapter{非线性振动初步}
\intro{本章介绍非线性振动的基本概念、几何理论与经典方程。非线性系统表现出叠加原理失效、频率-振幅耦合、跳跃现象、极限环和混沌等独特行为，揭示振动的丰富复杂性。}

\section{非线性振动的来源}
\subsection{几何非线性}
在小变形假设下，应变-位移关系可线性化。当位移幅度较大时，必须保留高阶项，导致运动方程中出现非线性恢复力项。例如，悬臂梁大挠度弯曲时，曲率表达式需使用精确形式：
\[
\kappa = \frac{w''}{[1+(w')^{2}]^{3/2}}
\]
展开后导出三次非线性项 $\frac{3}{2}(w')^{2}w''$。典型的几何非线性还包括：
\begin{itemize}
    \item 弦的大振幅振动（轴向伸长产生非线性张力）
    \item 薄板的大挠度（von K\'arm\'an 板理论）
    \item 柔性机构的几何非线性刚度
\end{itemize}

\subsection{材料非线性}
当应变超出材料的线弹性范围时，应力-应变关系偏离胡克定律，产生材料非线性：
\begin{itemize}
    \item \textbf{弹塑性}：屈服后应力-应变为非线性关系
    \item \textbf{超弹性}：橡胶等材料的 $S$ 形应力-应变曲线
    \item \textbf{黏弹性}：复模量随频率和温度变化
\end{itemize}
许多工程材料在小应变时表现为线弹性，在中大应变时呈现硬化或软化行为。

\subsection{接触与间隙非线性}
结构中存在间隙、碰撞、摩擦接触时，系统刚度发生突变，表现为分段线性或强非线性：
\begin{itemize}
    \item 齿轮齿侧间隙
    \item 裂纹的开合（呼吸裂纹）
    \item 间隙支撑的冲击振动
    \item 干摩擦阻尼接合面
\end{itemize}
这类系统的恢复力常表示为分段线性函数，导致复杂的动力学行为。

\subsection{非线性弹簧的力学特性}
\subsubsection{硬弹簧（Hardening Spring）}
恢复力随位移增长率超过线性比例：
\[
F_k(x) = k_1 x + k_3 x^{3}, \quad k_3 > 0
\]
随振幅增大，等效刚度增加，固有频率升高。典型的硬弹簧实例包括：
\begin{itemize}
    \item 两端固支梁的中等振幅振动（轴向拉伸刚化效应）
    \item 橡胶支座的大变形压缩
    \item 板壳在中大振幅下的膜效应
\end{itemize}

\begin{myformula}
硬弹簧的势能函数：
\[
V(x) = \int_0^{x} (k_1 \xi + k_3 \xi^{3})\,d\xi = \frac{1}{2}k_1 x^{2} + \frac{1}{4}k_3 x^{4}, \quad k_3 > 0
\]
势阱比线性弹簧的抛物形势阱更陡，大振幅时约束更强。
\end{myformula}

\subsubsection{软弹簧（Softening Spring）}
恢复力随位移增长低于线性比例（或 $k_3 < 0$）：
\[
F_k(x) = k_1 x + k_3 x^{3}, \quad k_3 < 0
\]
随振幅增大，等效刚度减小，固有频率降低。当 $|k_3|$ 足够大且振幅超过临界值时可出现静态不稳定（屈曲）。

\subsection{非线性阻尼}
阻尼力与速度的关系偏离线性关系时产生非线性阻尼：
\begin{itemize}
    \item \textbf{二次阻尼}：$F_d = c_2 |\dot{x}|\dot{x}$，流体在高雷诺数下的阻力
    \item \textbf{干摩擦（Coulomb 阻尼）}：$F_d = \mu N \sgn(\dot{x})$
    \item \textbf{黏弹性阻尼}：与频率和振幅相关的复刚度
    \item \textbf{材料迟滞阻尼}：加载-卸载循环中的能量耗散
\end{itemize}

\begin{figure}[H]
\centering
\begin{tikzpicture}[scale=1.0]
    \draw[-Stealth] (-3,0) -- (3,0) node[right] {$x$};
    \draw[-Stealth] (0,-2.5) -- (0,2.5) node[above] {$F_k(x)$};
    \draw[blue,thick] plot[domain=-2.5:2.5,samples=100] (\x,{0.6*\x + 0.15*\x*\x*\x});
    \draw[red,thick] plot[domain=-2.5:2.5,samples=100] (\x,{0.6*\x - 0.08*\x*\x*\x});
    \draw[gray,dashed] (-2.5,-1.5) -- (2.5,1.5);
    \node[blue] at (2.0,2.0) {硬弹簧 $k_3>0$};
    \node[red] at (2.0,-0.5) {软弹簧 $k_3<0$};
    \node[gray] at (2.0,1.5) {线性 $k_3=0$};
\end{tikzpicture}
\caption{硬弹簧与软弹簧的恢复力-位移曲线。硬弹簧（蓝色）恢复力增长快于线性，软弹簧（红色）增长慢于线性，虚线为线性弹簧参考。}
\label{fig:ch10_spring_force}
\end{figure}

\section{非线性系统的特征}
\subsection{叠加原理的失效}
线性系统满足叠加原理：若 $x_1(t)$ 和 $x_2(t)$ 为解，则任意线性组合 $\alpha x_1 + \beta x_2$ 也是解。非线性系统不满足此性质。

\begin{definition}[叠加原理失效]
设有非线性算子 $\mathcal{N}$ 满足 $\mathcal{N}(\alpha x_1 + \beta x_2) \neq \alpha\mathcal{N}(x_1) + \beta\mathcal{N}(x_2)$，则两个解的线性叠加不再满足原方程。例如对 Duffing 方程：
\[
\ddot{x} + x + \varepsilon x^{3} = 0
\]
若 $x_1$ 和 $x_2$ 分别为解，则 $\ddot{(x_1+x_2)} + (x_1+x_2) + \varepsilon(x_1+x_2)^{3} = \varepsilon(3x_1^{2}x_2 + 3x_1 x_2^{2}) \neq 0$，叠加不成立。
\end{definition}

\subsection{固有频率随振幅的变化}
非线性系统的等效刚度 $k_{\text{eq}}$ 与振幅 $A$ 相关，导致固有频率成为振幅的函数：
\[
\omega_n(A) \approx \omega_0 \left(1 + \frac{3}{8}\varepsilon A^{2}\right) \quad \text{(弱非线性Duffing)}
\]
其中 $\omega_0 = \sqrt{k_1/m}$ 为零振幅时的固有频率。
\begin{itemize}
    \item \textbf{硬弹簧} $\varepsilon > 0$：频率随振幅增大而增大，幅频曲线向右弯曲
    \item \textbf{软弹簧} $\varepsilon < 0$：频率随振幅增大而减小，幅频曲线向左弯曲
\end{itemize}

\subsection{跳跃现象（Jump Phenomenon）}
当激励频率缓慢扫过共振区时，响应振幅会出现突变——即跳跃现象。这是非线性系统幅频响应的多值性导致的。

\begin{figure}[H]
\centering
\begin{tikzpicture}[scale=1.1]
    \def\a{0.05}
    \draw[-Stealth] (0,0) -- (8,0) node[right] {$\omega/\omega_n$};
    \draw[-Stealth] (0,0) -- (0,5) node[above] {$A$（振幅）};
    \draw[blue,thick,smooth,variable=\t,domain=0.3:8,samples=200] plot({\t},{3.5/(sqrt((1-\t*\t)^2 + 4*\a*\a*\t*\t) + 0.0001*(\t*\t > 0.8 ? (\t-0.8)^3 : 0))});
    \draw[red,thick,-Stealth] (3.8,3.0) -- (5.5,3.0);
    \draw[red,thick,-Stealth] (5.8,0.9) -- (4.2,0.9);
    \node[red] at (7.0,2.6) {向上跳跃};
    \node[red] at (3.0,1.2) {向下跳跃};
    \draw[gray,thin,dashed] (3.5,0) -- (3.5,4.5);
    \draw[gray,thin,dashed] (5.5,0) -- (5.5,4.5);
    \node at (3.5,-0.4) {$\omega_1$};
    \node at (5.5,-0.4) {$\omega_2$};
    \node at (0.5,4.5) {硬弹簧 Duffing};
    \draw[green!50!black,dashed,thick] plot[domain=0.3:8,samples=200] (\x,{2.8/(sqrt((1-\x*\x)^2 + 4*\a*\a*\x*\x))});
    \node[green!50!black] at (7.5,1.0) {线性};
\end{tikzpicture}
\caption{Duffing 硬弹簧的幅频跳跃现象。频率从低到高扫频时（右侧箭头），振幅在 $\omega_1$ 处向上跳跃;从高到低扫频时（左侧箭头），振幅在 $\omega_2$ 处向下跳跃。虚线为对应的线性系统响应。}
\label{fig:ch10_jump}
\end{figure}

跳跃现象的发生机制：
\begin{enumerate}
    \item 扫频上行：振幅沿上部稳定分支缓慢增大，到达临界频率 $\omega_2$ 附近后失稳，振幅骤降至下部稳定分支
    \item 扫频下行：振幅沿下部稳定分支增大，到达 $\omega_1$ 时向上跳跃到上部稳定分支
    \item $\omega_1 < \omega < \omega_2$ 为多解区间，存在三个可能的稳态响应（两个稳定、一个不稳定）
\end{enumerate}

\subsection{次谐波与超谐波共振}
非线性系统可将能量在不同频率成分间传递。当激励频率 $\omega$ 接近固有频率 $\omega_n$ 的整数倍或分数倍时，可激发共振：

\begin{itemize}
    \item \textbf{超谐波共振（Superharmonic resonance）}：$\omega \approx \omega_n / n$，激励频率为固有频率的分数，例如 $\omega \approx \omega_n/3$ 时激发三次超谐波共振
    \item \textbf{次谐波共振（Subharmonic resonance）}：$\omega \approx m \omega_n$，激励频率为固有频率的整数倍，如 $\omega \approx 3\omega_n$
    \item \textbf{组合共振（Combination resonance）}：多频激励时，$\omega_1 \pm \omega_2 \approx \omega_n$
\end{itemize}

\subsection{自激振动与极限环}
自激振动（Self-excited vibration）指系统从恒定能源中汲取能量并将其转换为振动能量，维持稳态周期运动：
\begin{itemize}
    \item 颤振（flutter）：风速超过临界速度时薄翼的弯曲-扭转耦合振荡
    \item 干摩擦引起的黏滑振动：机床导轨爬行、制动噪声
    \item 车削颤振：刀具-工件系统的再生颤振
\end{itemize}

\begin{definition}[极限环]
极限环是相平面上的孤立闭轨，相邻轨线均渐近趋近（稳定极限环）或远离（不稳定极限环）该闭轨。极限环对应系统在时域中的自激稳态周期运动，其振幅和周期由系统参数决定，与初始条件无关。
\end{definition}

\section{相平面分析}
\subsection{相平面概念}
对于二阶自治系统：
\[
\ddot{x} + f(x, \dot{x}) = 0
\]
定义状态变量 $y = \dot{x}$，将二阶方程化为一阶微分方程组：
\[
\begin{cases}
\dot{x} = y \\
\dot{y} = -f(x, y)
\end{cases}
\]
以 $x$ 为横轴、$y = \dot{x}$ 为纵轴的平面称为\textbf{相平面}（Phase Plane）。系统的解在相平面上形成的曲线称为\textbf{相轨迹}。

\begin{myformula}
相轨迹的切线方向由速度场决定：
\[
\frac{dy}{dx} = \frac{\dot{y}}{\dot{x}} = -\frac{f(x,y)}{y}
\]
相轨迹不可相交（除奇点外），且满足唯一性（由 Picard-Lindel\"of 定理保证）。
\end{myformula}

\begin{itemize}
    \item 相轨迹顺时针环绕（若 $y>0$ 则 $x$ 增加，若 $y<0$ 则 $x$ 减少）
    \item 闭轨对应周期运动;螺旋线对应阻尼振动
    \item 奇点对应平衡位置（$\dot{x}=0$, $\dot{y}=0$）
\end{itemize}

\subsection{奇点分类与线性化分析}
在平衡点 $(x_e, 0)$ 附近，将非线性函数 $f(x,y)$ 线性化：
\[
f(x, y) \approx f(x_e, 0) + \frac{\partial f}{\partial x}\bigg|_{e}(x-x_e) + \frac{\partial f}{\partial y}\bigg|_{e} y = \omega_n^{2} \Delta x + 2\zeta \omega_n y
\]
线性化方程组为：
\[
\frac{d}{dt}\begin{pmatrix} \Delta x \\ y \end{pmatrix} = \begin{pmatrix} 0 & 1 \\ -\omega_n^{2} & -2\zeta\omega_n \end{pmatrix} \begin{pmatrix} \Delta x \\ y \end{pmatrix}
\]
由 Jacobi 矩阵的特征值 $\lambda_{1,2} = -\zeta\omega_n \pm \omega_n\sqrt{\zeta^{2}-1}$ 判断奇点类型：

\begin{table}[H]
\centering
\caption{奇点类型判据}
\begin{tabular}{ccc}
\toprule
阻尼比 $\zeta$ & 特征值性质 & 奇点类型 \\
\midrule
$\zeta = 0$ & 纯虚数 $\pm i\omega_n$ & \textbf{中心} \\
$0 < \zeta < 1$ & 共轭复数，实部负 & \textbf{稳定焦点} \\
$-1 < \zeta < 0$ & 共轭复数，实部正 & \textbf{不稳定焦点} \\
$\zeta > 1$ & 相异负实根 & \textbf{稳定结点} \\
$\zeta < -1$ & 相异正实根 & \textbf{不稳定结点} \\
鞍点工况 & 一正一负实根 & \textbf{鞍点} \\
\bottomrule
\end{tabular}
\label{tab:ch10_singular_pts}
\end{table}

\begin{definition}[中心 vs 焦点]
\textbf{中心}：特征值为纯虚数对，相轨迹为同心椭圆族，Lyapunov 稳定但非渐近稳定。\\
\textbf{焦点}：特征值为实部非零的共轭复数对，相轨迹为向奇点汇聚或发散的螺旋线。
\end{definition}

\begin{definition}[鞍点]
鞍点具有一维稳定流形（特征方向上的稳定子空间）和一维不稳定流形。相轨迹沿稳定流形方向趋近鞍点，沿不稳定流形方向远离鞍点。在非线性系统中，稳定流形与不稳定流形可能相交形成同宿轨道或异宿轨道。
\end{definition}

\begin{figure}[H]
\centering
\begin{tikzpicture}[scale=1.0]
    \def\titleA{中心 (Center)}
    \def\titleB{稳定焦点 (Stable Focus)}
    \def\titleC{鞍点 (Saddle)}
    \def\titleD{稳定结点 (Stable Node)}

    \begin{scope}
        \draw[-Stealth] (-1.8,0) -- (1.8,0);
        \draw[-Stealth] (0,-1.8) -- (0,1.8);
        \foreach \r in {0.3,0.7,1.1,1.5} {
            \draw[blue,thick] (0,0) ellipse ({\r} and {0.7*\r});
        }
        \node at (0,-2.2) {\titleA};
    \end{scope}

    \begin{scope}[xshift=4.5cm]
        \draw[-Stealth] (-1.8,0) -- (1.8,0);
        \draw[-Stealth] (0,-1.8) -- (0,1.8);
        \foreach \r in {0.3,0.7,1.1,1.5} {
            \draw[red,thick,domain=0:360,samples=80,variable=\t]
                plot({0.6*exp(-0.4*\t/360)*\r*cos(\t+50)},{0.6*exp(-0.4*\t/360)*\r*sin(\t+50)});
        }
        \node at (0,-2.2) {\titleB};
    \end{scope}

    \begin{scope}[xshift=9.0cm]
        \draw[-Stealth] (-1.8,0) -- (1.8,0);
        \draw[-Stealth] (0,-1.8) -- (0,1.8);
        \draw[green!50!black,thick] plot[domain=-1.5:1.5,samples=50] (\x,{0.7*\x/(1+abs(\x))});
        \draw[green!50!black,thick] plot[domain=-1.5:1.5,samples=50] (\x,{-0.7*\x/(1+abs(\x))});
        \draw[green!50!black,thick] plot[domain=0.1:1.8,samples=50] (\x,{sqrt(\x*\x-0.04)});
        \draw[green!50!black,thick] plot[domain=0.1:1.8,samples=50] (\x,{-sqrt(\x*\x-0.04)});
        \draw[green!50!black,thick] plot[domain=0.1:1.8,samples=50] ({-\x},{sqrt(\x*\x-0.04)});
        \draw[green!50!black,thick] plot[domain=0.1:1.8,samples=50] ({-\x},{-sqrt(\x*\x-0.04)});
        \node at (0,-2.2) {\titleC};
    \end{scope}

    \begin{scope}[xshift=13.5cm]
        \draw[-Stealth] (-1.8,0) -- (1.8,0);
        \draw[-Stealth] (0,-1.8) -- (0,1.8);
        \draw[orange,thick] (-1.5,-1.5) -- (0,0) -- (1.5,1.5);
        \draw[orange,thick] (-1.5,1.5) -- (0,0) -- (1.5,-1.5);
        \draw[orange,thick] (-1.0,-1.5) .. controls (-0.3,-0.8) and (0.3,-0.8) .. (1.0,-0.3);
        \draw[orange,thick] (-1.5,1.0) .. controls (-0.8,0.3) and (-0.8,-0.3) .. (-0.3,-1.0);
        \node at (0,-2.2) {\titleD};
    \end{scope}
\end{tikzpicture}
\caption{奇点类型的相平面图：中心（闭合椭圆轨道）、稳定焦点（渐近螺旋收敛）、鞍点（双曲线型）、稳定结点（直接收敛）。}
\label{fig:ch10_singularities}
\end{figure}

\subsection{极限环理论}
极限环是自激振动在相平面上的几何表征。考虑更一般的平面自治系统：
\[
\dot{\mathbf{x}} = \mathbf{F}(\mathbf{x}), \quad \mathbf{x} \in \R^{2}
\]

\begin{theorem}[Poincar\'e-Bendixson 定理]
设 $D$ 为平面内的有限区域，其内无奇点。若轨线恒在 $D$ 内而不离开，则 $\omega$ 极限集必为下列三者之一：
\begin{enumerate}
    \item 奇点
    \item 闭轨（周期解）
    \item 奇点与联结它们的轨线构成的连通集
\end{enumerate}
该定理是证明极限环存在性的核心工具。
\end{theorem}

\begin{proposition}[极限环的稳定性判据]
设极限环 $\Gamma$ 的参数方程为 $\mathbf{x} = \mathbf{p}(t)$，周期为 $T$。则极限环的稳定性由 Floquet 乘子决定。对于平面系统，可通过计算变分方程的单值矩阵判断：
\begin{itemize}
    \item 若特征乘子 $|\rho_1| < 1$（另一个 $|\rho_2|=1$），极限环渐近稳定
    \item 若任一特征乘子 $|\rho| > 1$，极限环不稳定
\end{itemize}
等价的判据是二维系统沿闭轨的积分（Poincar\'e 指标法）。
\end{proposition}

\begin{figure}[H]
\centering
\begin{tikzpicture}[scale=1.2]
    \draw[-Stealth] (-2.5,0) -- (2.5,0) node[right] {$x$};
    \draw[-Stealth] (0,-2.5) -- (0,2.5) node[above] {$\dot{x}$};

    \draw[red,very thick] (0,0) circle (1.5);
    \node[red] at (1.05,1.35) {$\Gamma$};

    \draw[blue,thick,-Stealth] (0.5,0) .. controls (1.2,0.6) and (1.2,1.4) .. (0.3,1.8);
    \draw[blue,thick,-Stealth] (0.8,0) .. controls (1.5,0.8) and (1.6,1.5) .. (0.2,2.0);
    \draw[blue,thick,-Stealth] (2.0,0.4) .. controls (1.8,1.8) and (0.8,1.8) .. (-0.8,1.8);
    \draw[blue,thick,-Stealth] (0.3,0) .. controls (0.8,0.4) and (1.0,1.0) .. (0.5,1.5);

    \draw[green!50!black,thick,-Stealth] (-0.3,-0.8) .. controls (-1.2,-0.4) and (-1.5,0) .. (-1.2,0.8);
    \draw[green!50!black,thick,-Stealth] (0.8,0.8) .. controls (2.0,0.4) and (2.5,0) .. (2.2,-0.6);

    \node[blue] at (2.8,1.8) {内部收敛};
    \node[green!50!black] at (2.8,-0.8) {外部收敛};
    \node[red] at (2.8,1.0) {稳定极限环};
\end{tikzpicture}
\caption{稳定极限环：内外部轨线均向环线趋近。Van der Pol 振子即为典型的有稳定极限环的系统。}
\label{fig:ch10_limit_cycle}
\end{figure}

\subsection{等倾线法}
等倾线是相平面上斜率为常数的点的轨迹，提供了一种手绘相轨迹的几何方法。平衡点的类型决定了全局相图结构。

\section{经典非线性方程}
\subsection{Duffing 方程}
Duffing 方程是描述具有非线性恢复力的受迫振子的经典模型：
\[
\ddot{x} + 2\zeta\dot{x} + \alpha x + \beta x^{3} = F\cos\omega t
\]
其中 $\beta > 0$ 为硬弹簧，$\beta < 0$ 为软弹簧，$F$ 为激励幅值，$\omega$ 为激励频率。

\begin{myformula}
Duffing 方程的一阶 Hamilton 形式（无阻尼自由振动 $\zeta=0$, $F=0$）：
\[
\dot{x} = y, \quad \dot{y} = -\alpha x - \beta x^{3}
\]
Hamilton 函数（总能量）为：
\[
H(x,y) = \frac{1}{2}y^{2} + \frac{1}{2}\alpha x^{2} + \frac{1}{4}\beta x^{4}
\]
当 $\beta>0$ 时，$H$ 正定，唯一平衡点为原点（中心）;当 $\beta<0$ 且 $\alpha<0$ 时出现双稳态。
\end{myformula}

Duffing 方程的核心特性：
\begin{enumerate}
    \item \textbf{主共振}（$\omega \approx \omega_n$）：幅频曲线弯曲，出现跳跃现象
    \item \textbf{次谐波共振}：当 $\omega \approx 3\omega_n$ 时激发 $1/3$ 阶次谐波
    \item \textbf{混沌运动}：在特定参数范围内，方程呈现确定性混沌
\end{enumerate}

\subsection{Van der Pol 方程}
Van der Pol 方程是自激振动的经典模型，由 Balthasar van der Pol 在 1920 年代提出以描述电子管振荡器：
\[
\ddot{x} - \mu(1-x^{2})\dot{x} + x = 0, \quad \mu > 0
\]

\begin{definition}[Van der Pol 振子的阻尼特征]
非线性阻尼项 $-\mu(1-x^{2})\dot{x}$ 具有特殊的振幅依赖性质：
\begin{itemize}
    \item 当 $|x| < 1$ 时，$1-x^{2} > 0$，阻尼为负（注入能量）
    \item 当 $|x| > 1$ 时，$1-x^{2} < 0$，阻尼为正（耗散能量）
\end{itemize}
系统在 $|x|=1$ 附近达到能量平衡，形成稳定的极限环。
\end{definition}

\begin{myformula}
对于小参数 $\mu \ll 1$，Van der Pol 方程的近似解为：
\[
x(t) \approx 2\cos t, \quad \dot{x}(t) \approx -2\sin t
\]
极限环近似为半径为 $2$ 的圆周。当 $\mu$ 很大时，极限环变为弛豫振荡（relaxation oscillation），由慢变段和快速跳跃段交替构成。
\end{myformula}

\begin{remark}
Van der Pol 方程与心脏跳动、神经元放电、电子振荡器等自激振荡现象密切相关。该方程是非线性动力学中研究最为深入的系统之一，其数学性质对整个自激振动理论具有示范意义。
\end{remark}

\subsection{单摆的大振幅振动}
单摆的运动方程：
\[
\ddot{\theta} + \frac{g}{L}\sin\theta = 0
\]
令 $\omega_0 = \sqrt{g/L}$。当 $\theta$ 不满足小角度假设时，精确方程含 $\sin\theta \approx \theta - \theta^{3}/6 + \dots$ 的高阶非线性项。

\begin{myformula}
单摆的精确周期（第一类椭圆积分）：
\[
T = 4\sqrt{\frac{L}{g}}\,K(k), \quad k = \sin\frac{\theta_{\max}}{2}
\]
其中 $K(k) = \int_0^{\pi/2} \frac{d\phi}{\sqrt{1 - k^{2}\sin^{2}\phi}}$ 为第一类完全椭圆积分。
级数展开：
\[
T = T_0\left[1 + \frac{1}{4}\sin^{2}\frac{\theta_{\max}}{2} + \frac{9}{64}\sin^{4}\frac{\theta_{\max}}{2} + \cdots\right], \quad T_0 = 2\pi\sqrt{\frac{L}{g}}
\]
\end{myformula}

\begin{figure}[H]
\centering
\begin{tikzpicture}[scale=1.1]
    \draw[-Stealth] (-3.5,0) -- (3.5,0) node[right] {$\theta$};
    \draw[-Stealth] (0,-2.5) -- (0,2.5) node[above] {$\dot{\theta}$};

    \draw[gray,thin] (-3.14,-2.5) -- (-3.14,2.5);
    \draw[gray,thin] (3.14,-2.5) -- (3.14,2.5);

    \foreach \E in {0.3,0.6,0.9,1.2,1.5,1.8,2.1} {
        \pgfmathsetmacro\vv{sqrt(2*\E + 2*cos(deg(0)))}
        \draw[blue,thick,smooth,variable=\t,domain=-180:180,samples=100]
            plot({deg(\t)/57.3},{sqrt(2*\E + 2*cos(deg(\t)))});
        \draw[blue,thick,smooth,variable=\t,domain=-180:180,samples=100]
            plot({deg(\t)/57.3},{-sqrt(2*\E + 2*cos(deg(\t)))});
    }

    \foreach \E in {2.2,2.6,3.0,3.5} {
        \draw[red,thick,smooth,variable=\t,domain=-250:250,samples=150]
            plot({deg(\t)/57.3},{sqrt(2*\E + 2*cos(deg(\t)))});
        \draw[red,thick,smooth,variable=\t,domain=-250:250,samples=150]
            plot({deg(\t)/57.3},{-sqrt(2*\E + 2*cos(deg(\t)))});
    }

    \node[blue] at (3.5,-0.8) {摆动（闭轨）};
    \node[red] at (3.5,-1.8) {旋转（开轨）};
    \node at (-3.14,2.7) {$-\pi$};
    \node at (3.14,2.7) {$\pi$};
\end{tikzpicture}
\caption{单摆的相轨迹图。低能量时（蓝色闭轨）为周期摆动，高能量（红色开轨）时为单向旋转。$\theta = \pm\pi$ 处为鞍点，其稳定流形与不稳定流形构成分界线（separatrix）。}
\label{fig:ch10_pendulum_phase}
\end{figure}

\section{近似解法}
\subsection{摄动法（Perturbation Method）}
当非线性项较小时（$0 < \varepsilon \ll 1$），可寻求解的小参数展开。考虑弱非线性 Duffing 方程：
\[
\ddot{x} + \omega_0^{2} x + \varepsilon x^{3} = 0
\]
假设解展为 $\varepsilon$ 的幂级数：
\[
x(t,\varepsilon) = x_0(t) + \varepsilon x_1(t) + \varepsilon^{2} x_2(t) + \cdots
\]

\begin{myformula}
代入方程并按 $\varepsilon$ 的幂次分离：
\begin{align*}
\mathcal{O}(1):&\quad \ddot{x}_0 + \omega_0^{2} x_0 = 0 \\
\mathcal{O}(\varepsilon):&\quad \ddot{x}_1 + \omega_0^{2} x_1 = -x_0^{3} \\
\mathcal{O}(\varepsilon^{2}):&\quad \ddot{x}_2 + \omega_0^{2} x_2 = -3x_0^{2} x_1
\end{align*}
\end{myformula}

零阶解为 $x_0 = A\cos(\omega_0 t + \phi)$。代入第一阶方程：
\[
\ddot{x}_1 + \omega_0^{2} x_1 = -A^{3}\cos^{3}(\omega_0 t + \phi) = -\frac{A^{3}}{4}\left[3\cos(\omega_0 t + \phi) + \cos(3\omega_0 t + 3\phi)\right]
\]
注意右端含 $\cos(\omega_0 t + \phi)$，与齐次解同频——这导致\textbf{长期项}（$t\cos\omega_0 t$ 形式的增长项），破坏解的有界性。直接展开的摄动法需引入附加条件消除长期项。

\subsection{Lindstedt-Poincar\'e 法}
为消除长期项，Lindstedt-Poincar\'e 法同时展开频率和解：
\[
\omega^{2} = \omega_0^{2} + \varepsilon\omega_1^{2} + \varepsilon^{2}\omega_2^{2} + \cdots
\]
引入新时间变量 $\tau = \omega t$，则 $\ddot{x} = \omega^{2} x''$，方程变为：
\[
\omega^{2} x'' + \omega_0^{2} x + \varepsilon x^{3} = 0
\]
将 $\omega^{2}$ 和 $x$ 展开后，选择各阶 $\omega_i^{2}$ 以消除长期项。

\begin{myformula}
Lindstedt-Poincar\'e 法给出的 Duffing 方程频率-振幅关系为：
\[
\omega = \omega_0\left[1 + \frac{3\varepsilon}{8\omega_0^{2}}A^{2} - \frac{15\varepsilon^{2}}{256\omega_0^{4}}A^{4} + \mathcal{O}(A^{6})\right]
\]
其中 $A$ 为基频振幅。
\end{myformula}

\subsection{平均法（Method of Averaging）}
考虑弱非线性方程：
\[
\ddot{x} + \omega_0^{2} x = \varepsilon f(x, \dot{x}, t)
\]
当 $\varepsilon = 0$ 时，通解为 $x = a\cos(\omega_0 t + \beta)$。令：
\[
x(t) = a(t)\cos[\omega_0 t + \beta(t)], \quad \dot{x}(t) = -a(t)\omega_0\sin[\omega_0 t + \beta(t)]
\]
附加约束 $\dot{a}\cos\psi - a\dot{\beta}\sin\psi = 0$（其中 $\psi = \omega_0 t + \beta$），将原方程化为振幅和相位的慢变方程：
\[
\[\begin{aligned}
\dot{a} &= -\frac{\varepsilon}{\omega_0} f(a\cos\psi, -a\omega_0\sin\psi, t)\sin\psi \\
\dot{\beta} &= -\frac{\varepsilon}{a\omega_0} f(a\cos\psi, -a\omega_0\sin\psi, t)\cos\psi
\end{aligned}\]
\]

\begin{myformula}
平均法的核心——在一周期内对右端取平均：
\[
\dot{\bar{a}} = -\frac{\varepsilon}{2\pi\omega_0}\int_0^{2\pi} f(a\cos\psi, -a\omega_0\sin\psi)\sin\psi\,d\psi
\]
\[
\dot{\bar{\beta}} = -\frac{\varepsilon}{2\pi a \omega_0}\int_0^{2\pi} f(a\cos\psi, -a\omega_0\sin\psi)\cos\psi\,d\psi
\]
平衡点 $\dot{\bar{a}} = 0$ 给出稳态振幅，稳定性由 Jacobi 矩阵特征值决定。
\end{myformula}

\subsection{谐波平衡法（Harmonic Balance Method）}
对于周期激励下的非线性系统：
\[
\ddot{x} + f(x, \dot{x}) = F\cos\omega t
\]
假设稳态解可表为截断的 Fourier 级数：
\[
x(t) = a_0 + \sum_{n=1}^{N} \left[ a_n\cos(n\omega t) + b_n\sin(n\omega t) \right]
\]
代入原方程后，令各谐波分量系数相等，得到关于 $\{a_n, b_n\}$ 的非线性代数方程组。

\begin{myformula}
对 Duffing 方程 $\ddot{x} + 2\zeta\dot{x} + \alpha x + \beta x^{3} = F\cos\omega t$，取单谐波近似 $x = A\cos(\omega t - \phi)$，代入得：
\[
\[\begin{aligned}
-\omega^{2}A\cos(\omega t - \phi) - 2\zeta\omega A\sin(\omega t - \phi) + \alpha A\cos(\omega t - \phi) & \\
+ \beta A^{3}\cos^{3}(\omega t - \phi) &= F\cos\omega t
\end{aligned}\]
\]
利用 $\cos^{3}\theta = \frac{3}{4}\cos\theta + \frac{1}{4}\cos3\theta$，忽略三次谐波（谐波平衡的截断），令基频分量相等：
\[
\left[(\alpha - \omega^{2})A + \frac{3}{4}\beta A^{3}\right]\cos(\omega t - \phi) - 2\zeta\omega A \sin(\omega t - \phi) = F\cos\omega t
\]
由此得频率-振幅方程：
\[
\left[(\alpha - \omega^{2})A + \frac{3}{4}\beta A^{3}\right]^{2} + (2\zeta\omega A)^{2} = F^{2}
\]
\end{myformula}

\section{混沌振动简介}
\subsection{确定性混沌的特征}
混沌（chaos）是确定性非线性动力系统表现出的对初值极端敏感的长期不可预测行为：
\begin{itemize}
    \item \textbf{对初值的敏感依赖}：相邻轨线以指数速率分离
    \item \textbf{拓扑传递性}：系统在相空间有限区域内的任意邻域随时间演化均稠密覆盖该区域
    \item \textbf{稠密的周期轨线}：任意混沌轨线的邻域内存在无穷多不稳定周期轨线
\end{itemize}

\begin{definition}[Lyapunov 指数]
设系统 $\dot{\mathbf{x}} = \mathbf{F}(\mathbf{x})$ 的两条相邻轨线初始距离为 $\delta \mathbf{x}_0$，演化后为 $\delta\mathbf{x}(t)$，定义：
\[
\lambda = \lim_{t\to\infty}\lim_{\|\delta\mathbf{x}_0\|\to 0} \frac{1}{t}\ln\frac{\|\delta\mathbf{x}(t)\|}{\|\delta\mathbf{x}_0\|}
\]
\begin{itemize}
    \item $\lambda < 0$：轨线收敛（规则运动）
    \item $\lambda = 0$：中性稳定（周期或准周期运动）
    \item $\lambda > 0$：轨线指数分离（混沌的基本判据）
\end{itemize}
$n$ 维系统存在 $n$ 个 Lyapunov 指数（Lyapunov 谱），最大 Lyapunov 指数 $\lambda_1$ 决定混沌存在性。
\end{definition}

\subsection{Duffing 方程的混沌}
受迫 Duffing 方程 $\ddot{x} + \delta\dot{x} + \alpha x + \beta x^{3} = \gamma\cos\omega t$ 在特定参数范围内呈现混沌运动。以 $\alpha = -1$, $\beta = 1$（双势阱型）为例，系统具有两个稳定的平衡点 $x = \pm 1$ 和一个鞍点 $x = 0$。

为探测混沌，常用本参数范围的典型值：
\begin{itemize}
    \item 阻尼参数 $\delta = 0.1 \sim 0.3$
    \item 激励幅值 $\gamma = 0.2 \sim 0.8$
    \item 频率 $\omega = 0.8 \sim 1.5$
\end{itemize}

\subsection{Poincar\'e 截面与分岔}
Poincar\'e 截面法：受迫周期系统的周期为 $T = 2\pi/\omega$，每周期采样一次（频闪采样），将连续动力系统降维为离散映射——Poincar\'e 映射 $\mathbf{P}: \R^{2} \to \R^{2}$。

\begin{itemize}
    \item \textbf{周期-$k$ 运动}：Poincar\'e 截面上出现 $k$ 个孤立点
    \item \textbf{准周期运动}：Poincar\'e 截面上为闭合曲线
    \item \textbf{混沌运动}：Poincar\'e 截面上呈分形点集（奇怪吸引子）
\end{itemize}

\begin{definition}[分岔（Bifurcation）]
当系统参数连续变化时，相图拓扑结构发生质的改变，称为分岔。常见类型：
\begin{itemize}
    \item \textbf{鞍结分岔}：稳定和不稳定不动点碰撞消失
    \item \textbf{Hopf 分岔}：焦点改变稳定性并产生极限环
    \item \textbf{倍周期分岔}：周期 $T$ 运动变为周期 $2T$ 运动，逐级通向混沌
\end{itemize}
控制参数 $\gamma$ 的扫变分岔图直观展示规则运动至混沌的过渡路径。
\end{definition}

\section{Python 代码实现}
\subsection{Duffing 方程的数值求解}
\begin{mycode}{Duffing 方程的 Runge-Kutta 数值解与相图}
\begin{lstlisting}[language=Python]
import numpy as np
from scipy.integrate import solve_ivp
import matplotlib.pyplot as plt

def duffing(t, y, delta, alpha, beta, gamma, omega):
    x, v = y
    dxdt = v
    dvdt = -delta * v - alpha * x - beta * x**3 + gamma * np.cos(omega * t)
    return [dxdt, dvdt]

delta, alpha, beta, gamma, omega = 0.25, -1.0, 1.0, 0.3, 1.0
y0 = [0.5, 0.0]
t_span = (0, 500)
t_eval = np.linspace(*t_span, 20000)

sol = solve_ivp(duffing, t_span, y0, args=(delta, alpha, beta, gamma, omega),
                t_eval=t_eval, method='RK45', rtol=1e-9)
x, v = sol.y

fig, (ax1, ax2) = plt.subplots(1, 2, figsize=(12, 5))
ax1.plot(sol.t[-2000:], x[-2000:], lw=0.6)
ax1.set_xlabel('$t$'); ax1.set_ylabel('$x(t)$')
ax1.set_title('Duffing方程时程')

ax2.plot(x[-4000:], v[-4000:], lw=0.3, color='blue')
ax2.set_xlabel('$x$'); ax2.set_ylabel('$\\dot{x}$')
ax2.set_title('相图（Poincaré截面采样后段）')
ax2.set_aspect('equal')
plt.tight_layout(); plt.show()
\end{lstlisting}
\end{mycode}

\subsection{分岔图}
\begin{mycode}{Duffing 方程分岔图（激励幅值扫变）}
\begin{lstlisting}[language=Python]
import numpy as np
from scipy.integrate import solve_ivp

def bifurcation_diagram(gamma_range, n_transient=300, n_sample=100):
    delta, alpha, beta, omega = 0.25, -1.0, 1.0, 1.0
    T = 2 * np.pi / omega
    bifur = []
    for gamma in gamma_range:
        sol = solve_ivp(duffing, [0, (n_transient + n_sample) * T],
                        [0.5, 0.0], args=(delta, alpha, beta, gamma, omega),
                        max_step=T/50, rtol=1e-9)
        x_vals = sol.y[0, -n_sample:]
        bifur.extend([(gamma, xv) for xv in x_vals])
    return np.array(bifur)

gamma_vals = np.linspace(0.1, 0.5, 400)
data = bifurcation_diagram(gamma_vals)

plt.figure(figsize=(10, 6))
plt.plot(data[:, 0], data[:, 1], 'k.', markersize=0.3, alpha=0.6)
plt.xlabel('$\\gamma$'); plt.ylabel('$x$ (Poincaré)')
plt.title('Duffing 方程分岔图')
plt.tight_layout(); plt.show()
\end{lstlisting}
\end{mycode}

\section{习题与思考}
\begin{enumerate}
    \item 推导硬弹簧 Duffing 方程的单谐波平衡法幅频关系，并分析跳跃现象存在的参数条件。
    \item 用 Lindstedt-Poincar\'e 法求无阻尼 Duffing 方程 $\ddot{x} + \omega_0^{2}x + \varepsilon x^{3} = 0$ 的频率-振幅关系至 $\mathcal{O}(\varepsilon^{2})$。
    \item 分析 Van der Pol 方程在 $\mu \ll 1$ 极限下极限环的振幅和周期，用平均法验证。
    \item 用相平面分析方法，定性绘制硬弹簧 Duffing 方程（$\alpha > 0, \beta > 0$）与双势阱 Duffing 方程（$\alpha < 0, \beta > 0$）的相图差异。
    \item 用 Poincar\'e 截面法判断 Duffing 方程在不同参数下的运动状态（周期、准周期、混沌）。
    \item 推导单摆大振幅周期的积分表达式，将椭圆积分展至 $\theta_{\max}^{4}$ 阶并评估小角度近似的误差（$\theta_{\max} = 30^{\circ}, 60^{\circ}, 90^{\circ}$）。
    \item 解释：为什么非线性系统的共振峰会出现"弯曲"（skewed resonance）？与线性系统相比，弯曲方向由什么决定？
    \item 编程实现 Duffing 方程的相图绘制，验证跳跃现象（扫频上下行），画出系统的分岔图并对主要分岔类型进行标注。
\end{enumerate}
